% =====================================================================
% Green IT 2026 (MICS2-63): Response to Reviewers
% Submitted together with the final paper, as a second PDF.
%
% Typically 1 to 2 pages. No minimum.
%
% Single column on purpose: this is a letter, quoting your reviewers and
% answering them. The double-column measure of the paper makes quoted
% blocks unreadable, so this is the one document of the four that sets
% `onecolumn`. Same class, so it looks like the rest of your submission.
%
% Build:  latexmk        (from the folder above this one, not from src/)
% =====================================================================
\documentclass[conference,onecolumn,a4paper]{IEEEtran}

\usepackage[T1]{fontenc}
\usepackage{url}
\usepackage{amsmath}

% ---------------------------------------------------------------------
% Guidance notes. Each one tells you what its section owes.
% Delete a note as you write its section, and set \guidancetrue below to
% \guidancefalse before you submit: the PDF you hand in carries no guidance.
\newif\ifguidance
\guidancetrue
\newcommand{\guidance}[1]{\ifguidance{\itshape\small #1\par\medskip}\fi}

\begin{document}

\title{Response to Reviewers: \textit{<title of your paper>}}

\author{\IEEEauthorblockN{<Your Full Name>}
\IEEEauthorblockA{Master in Information and Computer Sciences,
University of Luxembourg\\
<your.name>@student.uni.lu\\[2pt]
% Build stamp: the date this file was last compiled.
\scriptsize Compiled \today.}}

\maketitle


\guidance{Every italic note in this document is guidance.
Delete each one as you write its section, then set
\texttt{\textbackslash guidancefalse} in the preamble: the PDF you submit
carries no guidance.}

% ---------------------------------------------------------------------
\section*{Summary of Changes}

\guidance{One paragraph here, 300 words at most. What changed in the paper
between v1 and the final version, and why. A reader should be able to stop at
the end of this paragraph and know an overview of what you updated.}

% ---------------------------------------------------------------------
\section*{Response to Reviewer 1}

\guidance{One numbered block per comment in the review, in the shape below.
Every comment gets its own block, including the ones you changed nothing for.}

\begin{enumerate}

\item
\begin{quote}\itshape
<the reviewer's first comment, quoted verbatim>
\end{quote}

Modifications made.
\guidance{What you changed and where: the section, and the sentence or the
table. A reviewer should find it without re-reading the paper. If nothing
changed, write ``no change''.}

Argumentation.
\guidance{Why that change, or why that refusal. Disagreeing is allowed and
scores well when argued. Conceding every point scores the same as ignoring
them all.

Weigh a refusal. Repairing the defects your reviewers found is marked in the
paper itself, so a point you decline that the instructor judges well founded
stays there as a defect and is marked there. A refusal earns its marks when
the argument holds.

Argue with the point.}

\item
\begin{quote}\itshape
<the reviewer's second comment, quoted verbatim>
\end{quote}

Modifications made.

Argumentation.

\end{enumerate}

\guidance{And so on, one block per comment, to the end of the review.}

% ---------------------------------------------------------------------
\section*{Response to Reviewer 2}

\guidance{Same structure as above.}

\end{document}
